\documentclass[a4paper]{article}
\usepackage[margin=0.45in,top=0.38in,bottom=0.38in]{geometry}
\usepackage{multicol}
\usepackage{enumitem}
\usepackage{fancyhdr}
\usepackage{xcolor}
\usepackage{tabularx}
\pagestyle{fancy}\fancyhf{}\renewcommand{\headrulewidth}{0pt}\renewcommand{\footrulewidth}{0pt}
\setlength{\columnsep}{0.24in}
\setlength{\parindent}{0pt}\setlength{\parskip}{0pt}
\setlist[enumerate]{leftmargin=1.35em,label=\bfseries\arabic*.,labelsep=0.28em,itemsep=2pt,topsep=0pt,parsep=0pt,partopsep=0pt}
\newcommand{\examheader}{%
\begin{center}\textbf{\large SUMMER EXAMINATION 2026}\\[-1pt]
\textbf{\normalsize Artificial Intelligence}\\[-1pt]
\textbf{BSIT - 3rd Semester}\quad Time: 90 Minutes\quad Max Marks: 80\\[-1pt]
\textcolor{black!65}{\rule{\linewidth}{0.45pt}}\\[-2pt]
\begin{tabularx}{\linewidth}{@{}X X X@{}}Name: \hrulefill & Roll No.: \hrulefill & Date: \hrulefill\end{tabularx}\\[-1pt]
\textcolor{black!65}{\rule{\linewidth}{0.45pt}}\end{center}}
\begin{document}
\examheader
\raggedcolumns
\begin{multicols*}{3}
\fontsize{7}{8.2}\selectfont
\begin{enumerate}
\item AI is best described as the study of agents that perceive and act\newline\hspace*{0.6em}(A) intelligent behavior \quad (B) database normalization \newline\hspace*{0.6em}(C) file compression \quad (D) packet framing
\item The Turing test evaluates whether a machine can\newline\hspace*{0.6em}(A) converse indistinguishably from a human \quad (B) sort numbers \newline\hspace*{0.6em}(C) store facts \quad (D) encrypt files
\item Weak AI is designed for\newline\hspace*{0.6em}(A) a limited task \quad (B) every intellectual task \newline\hspace*{0.6em}(C) only robotics \quad (D) only games
\item A rational agent chooses an action that maximizes\newline\hspace*{0.6em}(A) expected performance \quad (B) memory usage \newline\hspace*{0.6em}(C) program length \quad (D) randomness
\item In PEAS, the E stands for\newline\hspace*{0.6em}(A) Environment \quad (B) Execution \newline\hspace*{0.6em}(C) Encoding \quad (D) Evaluation
\item A fully observable environment provides the agent with\newline\hspace*{0.6em}(A) complete relevant state information \quad (B) only rewards \newline\hspace*{0.6em}(C) no sensors \quad (D) past actions only
\item A reflex agent selects actions mainly from\newline\hspace*{0.6em}(A) the current percept \quad (B) a learned value function \newline\hspace*{0.6em}(C) a planning graph \quad (D) a proof tree
\item A model-based agent maintains\newline\hspace*{0.6em}(A) an internal state \quad (B) only a lookup table \newline\hspace*{0.6em}(C) a random seed \quad (D) a utility invoice
\item A utility-based agent uses utility to\newline\hspace*{0.6em}(A) rank desirable outcomes \quad (B) tokenize text \newline\hspace*{0.6em}(C) prune syntax trees \quad (D) count database rows
\item An episodic task environment has actions that\newline\hspace*{0.6em}(A) do not affect later episodes \quad (B) always require planning \newline\hspace*{0.6em}(C) must be continuous \quad (D) have no observations
\item Breadth-first search expands the\newline\hspace*{0.6em}(A) shallowest frontier node \quad (B) lowest heuristic node \newline\hspace*{0.6em}(C) deepest node \quad (D) random node
\item Depth-first search uses a frontier organized as a\newline\hspace*{0.6em}(A) stack \quad (B) queue \newline\hspace*{0.6em}(C) priority table \quad (D) hash ring
\item Uniform-cost search selects the frontier path with\newline\hspace*{0.6em}(A) lowest accumulated cost \quad (B) lowest depth \newline\hspace*{0.6em}(C) highest heuristic \quad (D) most children
\item A heuristic function estimates\newline\hspace*{0.6em}(A) cost from a state to a goal \quad (B) the exact training label \newline\hspace*{0.6em}(C) network bandwidth \quad (D) number of agents
\item A star search evaluates a node with\newline\hspace*{0.6em}(A) f(n)=g(n)+h(n) \quad (B) f(n)=g(n)-h(n) \newline\hspace*{0.6em}(C) f(n)=h(n)/g(n) \quad (D) f(n)=g(n)h(n) only
\item A heuristic is admissible when it\newline\hspace*{0.6em}(A) never overestimates true remaining cost \quad (B) always equals true cost \newline\hspace*{0.6em}(C) is negative \quad (D) ignores the goal
\item If every edge costs 1, a path of depth 4 has cost\newline\hspace*{0.6em}(A) 4 \quad (B) 1 \newline\hspace*{0.6em}(C) 0 \quad (D) 8
\item Hill climbing can fail at a\newline\hspace*{0.6em}(A) local maximum \quad (B) root node only \newline\hspace*{0.6em}(C) terminal compiler \quad (D) balanced database
\item Simulated annealing sometimes accepts\newline\hspace*{0.6em}(A) a worse move \quad (B) only a goal move \newline\hspace*{0.6em}(C) no move \quad (D) a duplicate state
\item Minimax assumes that the opponent\newline\hspace*{0.6em}(A) acts optimally \quad (B) acts randomly \newline\hspace*{0.6em}(C) never moves \quad (D) maximizes our score
\item Alpha-beta pruning changes minimax results\newline\hspace*{0.6em}(A) never \quad (B) only at even depth \newline\hspace*{0.6em}(C) always \quad (D) only with chance nodes
\item In alpha-beta pruning, alpha is the best value found for\newline\hspace*{0.6em}(A) MAX so far \quad (B) MIN so far \newline\hspace*{0.6em}(C) the root heuristic \quad (D) a random leaf
\item A constraint satisfaction problem consists of variables, domains, and\newline\hspace*{0.6em}(A) constraints \quad (B) rewards \newline\hspace*{0.6em}(C) gradients \quad (D) tokens
\item Backtracking search usually assigns\newline\hspace*{0.6em}(A) one variable at a time \quad (B) all variables randomly \newline\hspace*{0.6em}(C) only goal variables \quad (D) no values
\item The MRV heuristic chooses the variable with\newline\hspace*{0.6em}(A) fewest legal values \quad (B) largest domain \newline\hspace*{0.6em}(C) highest cost \quad (D) most neighbors only
\item Arc consistency removes values that\newline\hspace*{0.6em}(A) have no supporting neighbor value \quad (B) are globally optimal \newline\hspace*{0.6em}(C) are duplicated \quad (D) are outside the root
\item In STRIPS, an action has preconditions and\newline\hspace*{0.6em}(A) effects \quad (B) pixels \newline\hspace*{0.6em}(C) gradients \quad (D) priors
\item Forward planning starts at\newline\hspace*{0.6em}(A) the initial state \quad (B) the goal state \newline\hspace*{0.6em}(C) the deepest leaf \quad (D) the utility maximum
\item Backward planning works backward from\newline\hspace*{0.6em}(A) the goal \quad (B) the first percept \newline\hspace*{0.6em}(C) the largest domain \quad (D) the reward table
\item A proposition is a statement that is\newline\hspace*{0.6em}(A) true or false \quad (B) always numerical \newline\hspace*{0.6em}(C) a probability distribution \quad (D) a neural layer
\item Forward chaining applies rules from\newline\hspace*{0.6em}(A) known facts toward conclusions \quad (B) goals toward facts \newline\hspace*{0.6em}(C) random facts \quad (D) only contradictions
\item Backward chaining begins with\newline\hspace*{0.6em}(A) a query goal \quad (B) an empty database \newline\hspace*{0.6em}(C) a sensor reading \quad (D) a heuristic value
\item A Bayesian network represents dependencies using a\newline\hspace*{0.6em}(A) directed acyclic graph \quad (B) stack \newline\hspace*{0.6em}(C) confusion matrix \quad (D) parse tree only
\item Naive Bayes assumes features are conditionally independent given the\newline\hspace*{0.6em}(A) class \quad (B) learning rate \newline\hspace*{0.6em}(C) action \quad (D) cluster index
\item Fuzzy logic allows truth values\newline\hspace*{0.6em}(A) between 0 and 1 \quad (B) only -1 and 1 \newline\hspace*{0.6em}(C) only 0 or 1 \quad (D) above 100
\item Supervised learning uses\newline\hspace*{0.6em}(A) labeled examples \quad (B) no observations \newline\hspace*{0.6em}(C) only rewards \quad (D) only rules
\item Unsupervised learning seeks structure in\newline\hspace*{0.6em}(A) unlabeled data \quad (B) test labels \newline\hspace*{0.6em}(C) utility functions \quad (D) action policies
\item Reinforcement learning learns from\newline\hspace*{0.6em}(A) rewards and penalties \quad (B) class labels only \newline\hspace*{0.6em}(C) SQL joins \quad (D) grammar rules
\item A validation set is used mainly for\newline\hspace*{0.6em}(A) model selection and tuning \quad (B) final unbiased reporting \newline\hspace*{0.6em}(C) data collection \quad (D) encryption
\item Overfitting means a model fits training data well but\newline\hspace*{0.6em}(A) generalizes poorly \quad (B) has no parameters \newline\hspace*{0.6em}(C) cannot train \quad (D) has low variance always
\end{enumerate}
\end{multicols*}
\newpage
\examheader
\raggedcolumns
\begin{multicols*}{3}
\fontsize{7}{8.2}\selectfont
\begin{enumerate}[resume]
\item High bias commonly causes\newline\hspace*{0.6em}(A) underfitting \quad (B) memorization \newline\hspace*{0.6em}(C) data leakage \quad (D) high variance only
\item Precision is the fraction of predicted positives that are\newline\hspace*{0.6em}(A) actually positive \quad (B) actually negative \newline\hspace*{0.6em}(C) all observations \quad (D) missed positives
\item Recall is the fraction of actual positives that are\newline\hspace*{0.6em}(A) correctly found \quad (B) predicted negative \newline\hspace*{0.6em}(C) all negatives \quad (D) true negatives only
\item For TP=8 and FP=2, precision is\newline\hspace*{0.6em}(A) 0.80 \quad (B) 0.20 \newline\hspace*{0.6em}(C) 0.60 \quad (D) 0.90
\item The F1 score is the harmonic mean of\newline\hspace*{0.6em}(A) precision and recall \quad (B) accuracy and loss \newline\hspace*{0.6em}(C) bias and variance \quad (D) speed and memory
\item A perceptron computes a weighted sum followed by\newline\hspace*{0.6em}(A) an activation \quad (B) a database join \newline\hspace*{0.6em}(C) a graph search \quad (D) a token split
\item Gradient descent updates parameters using the\newline\hspace*{0.6em}(A) negative loss gradient \quad (B) positive labels only \newline\hspace*{0.6em}(C) largest feature \quad (D) random action
\item A convolutional neural network is especially suited to\newline\hspace*{0.6em}(A) spatial patterns \quad (B) symbolic proofs only \newline\hspace*{0.6em}(C) transaction logs only \quad (D) shortest paths
\item Pooling in a CNN commonly reduces\newline\hspace*{0.6em}(A) spatial dimensions \quad (B) the number of classes \newline\hspace*{0.6em}(C) training labels \quad (D) token vocabulary only
\item An RNN is designed to process\newline\hspace*{0.6em}(A) sequences \quad (B) static tables only \newline\hspace*{0.6em}(C) graphs only \quad (D) images without order
\item An LSTM uses gates to help preserve\newline\hspace*{0.6em}(A) long-range dependencies \quad (B) database keys \newline\hspace*{0.6em}(C) pixel colors only \quad (D) search depth
\item A Markov decision process includes states, actions, transition model, rewards, and\newline\hspace*{0.6em}(A) a discount factor \quad (B) a compiler \newline\hspace*{0.6em}(C) a parse tree \quad (D) a domain vocabulary
\item A policy maps states to\newline\hspace*{0.6em}(A) actions or action probabilities \quad (B) class labels only \newline\hspace*{0.6em}(C) database rows \quad (D) heuristic costs
\item Q-learning estimates\newline\hspace*{0.6em}(A) action values \quad (B) class priors only \newline\hspace*{0.6em}(C) parse probabilities \quad (D) pixel edges
\item The Q-learning update uses the reward plus discounted\newline\hspace*{0.6em}(A) next-state best Q value \quad (B) current input size \newline\hspace*{0.6em}(C) training accuracy \quad (D) heuristic depth
\item Exploration in reinforcement learning means\newline\hspace*{0.6em}(A) trying uncertain actions \quad (B) always choosing the best known action \newline\hspace*{0.6em}(C) removing rewards \quad (D) stopping training
\item A false positive occurs when a negative case is\newline\hspace*{0.6em}(A) predicted positive \quad (B) predicted negative \newline\hspace*{0.6em}(C) left unlabeled \quad (D) counted twice
\item Data leakage occurs when training uses information\newline\hspace*{0.6em}(A) unavailable at prediction time \quad (B) from a larger batch \newline\hspace*{0.6em}(C) with missing values \quad (D) after normalization
\item Responsible AI emphasizes fairness, transparency, privacy, and\newline\hspace*{0.6em}(A) accountability \quad (B) unlimited surveillance \newline\hspace*{0.6em}(C) opaque decisions \quad (D) forced automation
\item A model card is intended to document\newline\hspace*{0.6em}(A) model purpose, limits, and performance \quad (B) network routes \newline\hspace*{0.6em}(C) database indexes \quad (D) source code licenses only
\item Differential privacy protects individuals by adding controlled\newline\hspace*{0.6em}(A) noise \quad (B) labels \newline\hspace*{0.6em}(C) search depth \quad (D) network hops
\item A vector dot product measures alignment and is largest for vectors pointing\newline\hspace*{0.6em}(A) in similar directions \quad (B) in opposite directions \newline\hspace*{0.6em}(C) at random \quad (D) with zero length only
\item The determinant of a 2 by 2 matrix [[a,b],[c,d]] is\newline\hspace*{0.6em}(A) ad-bc \quad (B) ab+cd \newline\hspace*{0.6em}(C) ac-bd \quad (D) a+b+c+d
\item If a fair binary event has probability 0.5, its odds of success are\newline\hspace*{0.6em}(A) 1 to 1 \quad (B) 1 to 2 \newline\hspace*{0.6em}(C) 2 to 1 \quad (D) 0 to 1
\item Cross-entropy penalizes a confident prediction that is\newline\hspace*{0.6em}(A) wrong \quad (B) correct \newline\hspace*{0.6em}(C) uniform \quad (D) missing a feature
\item A confusion matrix summarizes\newline\hspace*{0.6em}(A) classification outcomes \quad (B) regression slopes \newline\hspace*{0.6em}(C) search paths \quad (D) cluster centroids
\item A ROC curve plots true positive rate against\newline\hspace*{0.6em}(A) false positive rate \quad (B) true negative rate \newline\hspace*{0.6em}(C) loss \quad (D) sample size
\item K-means assigns points to the nearest\newline\hspace*{0.6em}(A) centroid \quad (B) decision stump \newline\hspace*{0.6em}(C) support vector \quad (D) parse node
\item PCA projects data onto directions of greatest\newline\hspace*{0.6em}(A) variance \quad (B) class count \newline\hspace*{0.6em}(C) reward \quad (D) text frequency only
\item The principal purpose of normalization in ML is to\newline\hspace*{0.6em}(A) put features on comparable scales \quad (B) create labels \newline\hspace*{0.6em}(C) increase leakage \quad (D) remove all outliers
\item A learning rate that is too large may cause training to\newline\hspace*{0.6em}(A) diverge \quad (B) stop overfitting automatically \newline\hspace*{0.6em}(C) produce labels \quad (D) become deterministic
\item A learning rate that is too small usually makes convergence\newline\hspace*{0.6em}(A) slow \quad (B) impossible in every case \newline\hspace*{0.6em}(C) random \quad (D) independent of data
\item An agent's sensors provide percepts, while its actuators produce\newline\hspace*{0.6em}(A) actions \quad (B) heuristics \newline\hspace*{0.6em}(C) constraints \quad (D) utilities
\item Iterative deepening combines the low memory of DFS with the completeness of\newline\hspace*{0.6em}(A) BFS \quad (B) hill climbing \newline\hspace*{0.6em}(C) minimax \quad (D) beam search
\item A goal test checks whether a state satisfies\newline\hspace*{0.6em}(A) the goal condition \quad (B) the heuristic estimate \newline\hspace*{0.6em}(C) the action cost \quad (D) the percept history
\item Knowledge engineering converts domain expertise into\newline\hspace*{0.6em}(A) a knowledge base \quad (B) a pixel matrix \newline\hspace*{0.6em}(C) a routing table \quad (D) a loss curve
\item The utility of an outcome represents its\newline\hspace*{0.6em}(A) desirability \quad (B) token count \newline\hspace*{0.6em}(C) search depth \quad (D) probability only
\item An expert system usually contains a knowledge base and an\newline\hspace*{0.6em}(A) inference engine \quad (B) image encoder \newline\hspace*{0.6em}(C) order queue \quad (D) optimizer
\item The discount factor gamma is normally constrained to the interval\newline\hspace*{0.6em}(A) 0 to 1 \quad (B) 1 to 2 \newline\hspace*{0.6em}(C) -2 to -1 \quad (D) 10 to 100
\item An AI agent's performance measure specifies\newline\hspace*{0.6em}(A) how success is evaluated \quad (B) how files are compressed \newline\hspace*{0.6em}(C) how tokens are split \quad (D) how keys are stored
\end{enumerate}
\end{multicols*}
\end{document}
